% Generated draft: requires mathematical and source review.
\documentclass[11pt]{article}
\usepackage[margin=25mm]{geometry}
\usepackage{fontspec}
\setmainfont{Latin Modern Roman}
\usepackage{amsmath,amssymb,mathtools,unicode-math}
\setmathfont{Latin Modern Math}
\usepackage{enumitem,needspace,xurl,hyperref}
\setlength{\parindent}{0pt}
\setlength{\parskip}{5pt}
\setlength{\emergencystretch}{3em}
\newcommand{\N}{\mathbb{N}}
\newcommand{\Z}{\mathbb{Z}}
\newcommand{\Q}{\mathbb{Q}}
\newcommand{\R}{\mathbb{R}}
\newcommand{\C}{\mathbb{C}}
\newcommand{\F}{\mathbb{F}}
\newcommand{\A}{\mathbb{A}}
\newcommand{\PP}{\mathbb{P}}
\newcommand{\E}{\mathbb{E}}
\newcommand{\eps}{\varepsilon}
\newcommand{\dd}{\,\mathrm{d}}
\begin{document}
\section*{UnsolvedMath ID: 800403}
\textbf{Source problem code:} catalogue entry
\subsection{Definitions and mathematical statement}
Let \(K\) denote the ground field of the projective space in the source; the supplied record does not specify \(K\), its characteristic, or any restrictions on it. In the usual positive-dimensional setting, let \(N\geq 1\) and let
\[
f:\PP^N_K\dashrightarrow\PP^N_K
\]
be a dominant rational map, meaning that its image is Zariski dense. For \(n\geq 1\), let \(f^n\) denote its \(n\)-fold iterate. The algebraic degree \(\deg(f^n)\) is the common degree of homogeneous coordinate polynomials representing \(f^n\), after their common factors have been removed. Interpreting the source's dynamical degree as the first dynamical degree, set
\[
\delta_f=\lim_{n\to\infty}\bigl(\deg(f^n)\bigr)^{1/n}.
\]
The supplied \emph{Integrality Conjecture}, attributed to Bellon--Viallet (1999), asserts that for every such map, \(\delta_f\) is an algebraic integer: there exist an integer \(m\geq 1\) and coefficients \(a_0,\ldots,a_{m-1}\in\Z\) such that
\[
\delta_f^m+a_{m-1}\delta_f^{m-1}+\cdots+a_1\delta_f+a_0=0.
\]
The intended ground-field hypotheses remain unresolved in the supplied excerpt. The positive-dimensional convention and the interpretation as the first dynamical degree require confirmation against the source context.
\subsection{Short English statement}
The conjecture says that the exponential growth rate of the algebraic degrees of iterates of every dominant rational self-map of projective space is an algebraic integer. The supplied excerpt leaves the ground field unspecified. Its registry status is \emph{open}; this is a supplied classification, not an independently verified assessment.
\subsection{Sources}
\begin{itemize}
\item Local UnsolvedMath registry: permanent ID \(800403\), exact source problem code \texttt{catalogue entry}; title: \emph{Open Problems in Arithmetic Dynamics --- Conjecture 26 (Integrality Conjecture)}. The supplied record labels the problem open.
\item Benedetto, DeMarco, Ingram, Jones et al., \emph{Current Trends and Open Problems in Arithmetic Dynamics} (2018), Conjecture 26. Supplied source URL: \url{https://arxiv.org/abs/1806.04980}. The record reports extraction by \texttt{arxiv\_latex} and access on 2026-10-16; these are supplied provenance details, not an independently verified visit.
\item Original conjecture attribution in the supplied statement: Bellon--Viallet (1999), reference label \texttt{MR1704282}. No further bibliographic details for that reference are supplied. The registry's proposer field lists Benedetto, DeMarco, Ingram, Jones et al., with year 2018, corresponding to the supplied survey attribution.
\end{itemize}
\end{document}
